\documentclass[11pt,letterpaper]{article}
\usepackage[margin=0.8in]{geometry}
\usepackage[T1]{fontenc}
\usepackage{lmodern}
\usepackage{amsmath,amssymb,xcolor,listings,enumitem,hyperref}
\hypersetup{colorlinks=true,linkcolor=blue,urlcolor=blue,plainpages=false,pdftitle={CMSC433 Topic Practice}}
\definecolor{paper}{RGB}{246,248,250}
\lstset{basicstyle=\ttfamily\small,backgroundcolor=\color{paper},columns=fullflexible,
keepspaces=true,breaklines=true,showstringspaces=false,frame=single,
rulecolor=\color{gray!30},aboveskip=0.5em,belowskip=0.5em}
\lstnewenvironment{codeanswer}{\lstset{backgroundcolor=\color{white},frame=single}}{}
\setlist[enumerate]{itemsep=0.7em}
\setlength{\parindent}{0pt}
\setlength{\parskip}{0.5em}
\ExplSyntaxOn
\NewDocumentCommand{\code}{m}{\texttt{%
\tl_set:Nx\l_tmpa_tl{\tl_to_str:n{#1}}%
\seq_set_split:NnV\l_tmpa_seq{~}\l_tmpa_tl
\seq_map_indexed_inline:Nn\l_tmpa_seq{%
\int_compare:nNnT{##1}>{1}{\penalty0\relax\space}%
\str_map_inline:nn{##2}{####1\kern0pt\relax}}}}
\ExplSyntaxOff
% TYPE YOUR WRITTEN ANSWERS INSIDE THESE BRACES in the included lesson.tex files.
\newcommand{\answer}[1]{\par\nopagebreak[4]\smallskip\noindent\textbf{Answer: }#1\par\medskip}
\newcommand{\task}[1]{\subsection*{#1}}
\newcommand{\source}[1]{{\small\textbf{Source:} #1}\par}
\newcommand{\ready}[1]{{\small\textbf{Checkpoint:} #1}\par}

\raggedbottom
\newcommand{\sheettitle}[3]{\noindent{\small CMSC433 / Part #1}\par
{\Large\bfseries #2}\par{\large #3}\par\smallskip\hrule\medskip}
\newcommand{\mockrules}{{\small Topic mock, using the sample's question styles.
Attempt independently; feedback after the whole file. Untimed practice, no points
assigned. Current exam scope and rules are unconfirmed.}\par}

\begin{document}
\sheettitle{01}{Type foundations}{Guided practice}

\textbf{Use with the chat:} read the idea and example, then do Question 1 only.
We discuss each attempt before continuing.

\textbf{Idea.} Application matches an argument's type with a function's input.
Keep lists, pairs, and constraints in your answer. Arrows group to the right;
function application groups to the left.

\textbf{Worked example.} In \code{\h x -> h (Just x)}, let \code{x :: a}.
Then \code{Just x :: Maybe a}, so \code{h :: Maybe a -> b}.
The whole type is \code{(Maybe a -> b) -> a -> b}.

\task{1. Match the shapes}
Infer \code{\g x -> g [x]}. Fill the ledger before writing the full type.
\answer{x :: [ ]; [x] :: [ ]; g :: [ ]; result :: [ ]; whole lambda :: [ ].}

\task{2. Keep the necessary constraint}
Give the most general type of \code{\xs -> [even x | x <- xs]}.
Show which operation introduces a constraint; do not default to \code{Int}.
\answer{[Type and reasoning.]}

\task{3. Diagnose a conflict}
Give the most general type or explain why \code{\f -> (f True, f 'q')}
is ill-typed. State the incompatible requirements on \code{f}.
\answer{[Type or conflicting equations.]}

\source{Sample Q1, p. 1; Monads/GenericMonads lecture signatures.}

\end{document}
